Đồ án tốt nghiệp này gồm 6 chương với cấu trúc như sau.
\begin{itemize}
    \item \textbf{Chương 1: Tổng quan}\\
Chương này sẽ để cập tổng quan về đề tài của đồ án, cùng với động cơ, mục đích tiến hành đề tài và lợi ích mà đề tài có thể mang lại; sau đó, trình bày về mục tiêu thực hiện, phạm vi, ý nghĩa và cấu trúc của dồ án.
    \item \textbf{Chương 2: Cơ sở lí thuyết}\\
Chương này sẽ đề cập đến những kiến thức nền tảng cần chuẩn bị để thực hiện được đồ án này, bao gồm kiến thức cho các lựa chọn để tiếp cận dự án và các kiến thức liên quan đến các công nghệ được sử dụng cho dự án.
    \item \textbf{Chương 3: Các công trình liên quan}
    \begin{itemize}
        \item[$\circ$] Tổng quan về các dự án, hệ thống kết nối trong lĩnh vực AI hoặc công nghệ nói chung có liên quan tới AI.
        \item[$\circ$] Phân tích ưu nhược điểm của các hệ thống này.
    \end{itemize}
    \item \textbf{Chương 4: Phân tích hệ thống}\\
Chương này sẽ trình bày về quá trình phân tích hệ thống, bao gồm phân tích yêu cầu, phân tích chức năng, phân tích luồng dữ liệu và phân tích các vấn đề,rủi ro có thể xảy ra và đưa ra các hướng giải quyết phù hợp. Qua đó, dự án sẽ có một cơ sở tốt để tiến hành thiết kế hệ thống.
    \item \textbf{Chương 5: Thiết kế hệ thống}\\
Chương này sẽ trình bày về quá trình thiết kế hệ thống, bao gồm thiết kế kiến trúc, thiết kế cơ sở dữ liệu, thiết kế giao diện người dùng và thiết kế các chức năng cụ thể. Qua đó, đồ án sẽ có được một hệ thống hoàn chỉnh, phù hợp với yêu cầu và mục đích của dự án.
    \item \textbf{Chương 6: Kiểm thử, triển khai và đánh giá hệ thống}\\
    Kiểm thử hệ thống, triển khai hệ thống lên server và đánh giá hiệu năng của hệ thống
    \item \textbf{Chương 7: Kết luận}
    \begin{itemize}
        \item[$\circ$] Tóm tắt lại nội dung và kết quả đạt được.
        \item[$\circ$] Đánh giá tổng quan về đề tài.
        \item[$\circ$] Đề xuất hướng phát triển và công trình tương lai.
    \end{itemize}
    .
\end{itemize}

